\section{Research path and conclusion}
\label{sec:path}

The program should advance in an order in which each step can fail.

\paragraph{Role succession in a living system.} Carry out the prospective
\textit{INO1} test of Section~\ref{sec:ino1}, and its analogues in other
systems with experimentally separable record and role: prion conformational
inheritance, CRISPR spacer acquisition and loss, and enzyme networks with
measured turnover. Freeze the role equivalence, horizon, and shams before
opening outcomes, and fit the best ordinary model on training interventions
before evaluating both on held-out ones.

\paragraph{Replication across artificial laws.} Repeat the testbed's
succession protocols under alternative local laws, schedulers, topologies, and
independent implementations, requiring
Eq.~\eqref{eq:icp-constraint-divergence} to keep its qualitative result.

\paragraph{Adjacent-frontier assay.} On each system, freeze the type map,
actualized repertoire, policies, and horizons; report enabled, lost, and
retained identities (Eq.~\eqref{eq:possibility-flux-ledger}) rather than novelty
counts; and compare against fixed-state, reinforced-urn, branching, and TAP
rivals. Test the carrier dependence of the frontier
(Section~\ref{sec:ap}) by removing and restoring a carrier after its novelty
event, and attempt to predict a system's triggering rate from carrier
measurements before fitting its sequence. The strongest near-term signature is a prospectively predicted
crossover of the form of Eq.~\eqref{eq:constraint-enablement-crossover}, with
$h_\star$ forecast from separately measured carrier lifetimes.

\paragraph{Stopping rules.} The principle should be narrowed or reformulated
if complete-state augmentation repeatedly removes the candidate historical
variable, if carrier interventions fail under matched controls, if the role
disappears under reasonable representation changes, or if no compression,
invariant, prohibition, or prediction emerges beyond native theories. A result
that reduces ICP to a precise theory of physically embodied memory would still
be valuable.

\paragraph{Conclusion.} Information can persist beyond the object that first
carried it, but every established example does so through continuing physical
realization in another state, relation, or organization. Biology shows records
outliving their initiators, constraints maintained through the turnover of
their realizers, and futures opened by prior organization. ICP's contribution
is a single operational test, with a signed formalism for pruning and enabling,
that makes these cases comparable and that separates the persistence of a
record from the persistence of a role. Its claim to be more than an organizing
vocabulary rests on a result it has not yet produced: a held-out prediction,
from carrier succession, that the best complete-state model of the same system
misses.
